Los resultados confirman que la complejidad asintótica no basta para
predecir el desempeño práctico. \textsc{Strassen}, pese a su menor
complejidad teórica ($O(n^{2.807})$ frente a $O(n^3)$ de \textsc{Naive}),
resultó hasta $364\times$ más lento en el rango medido, debido al
overhead de asignación de memoria de su recursión.
De forma análoga, la optimización de mediana de tres en \textsc{QuickSort},
diseñada para evitar el peor caso clásico del pivote fijo, introdujo una
vulnerabilidad distinta (el fenómeno conocido como \textit{"median-of-three
killer"} \cite{mcilroy1999}): en el dominio D7, tardó $331\times$ y
$265\times$ más en los casos ascendente y descendente respectivamente
que en D1, mientras que en el caso aleatorio la diferencia fue de solo
$7.4\times$. En ambos experimentos se repite el mismo patrón: una
optimización teóricamente sólida puede fallar en la práctica bajo
condiciones de entrada específicas que la notación asintótica no
captura, reforzando la necesidad de validación experimental más allá
del análisis teórico.